% ===========================================================================
%  Appendice E — Inventario delle run
% ===========================================================================
\chapter{Inventario delle run}
\label{app:inventario}

Questa appendice elenca ogni esecuzione utilizzata nella tesi. Serve a rendere
verificabile l'affermazione, posta nella \cref{sec:minacce}, che nessuna run
storica a schema diverso sia stata mescolata alle run della campagna definitiva.

\begin{center}\small
\begin{tabular}{@{}l l l l@{}}\toprule
\textbf{Esp.}&\textbf{Unità appaiate}&\textbf{Artefatto canonico}&\textbf{Stato}\\\midrule
01&10 seed, 20 run&\code{experiment\_01\_quantum\_gpt\_n10.json}&completo\\
02&5 seed, 10 run&\code{experiment\_02\_quantum\_vit\_n5.json}&completo\\
03&3 dataset, 5 seed, 45 run&\code{experiment\_03\_quantum\_reg\_n100\_steps400.json}&completo\\
04&20 sottocampioni&\code{experiment\_04\_quantum\_kernel\_frozen.json}&completo\\\bottomrule
\end{tabular}\end{center}

Gli artefatti sono in \file{docs/checkpoints}; l'elenco degli hash SHA-256 è
in \file{docs/HANDOFF.md}. I manifest originali registrano il commit
\code{b91c318ce87e697d7da459cd209a4c11a742fb11} e \code{dirty=true}. I percorsi
delle singole run, seed, tempi, configurazioni e versioni di schema sono
preservati negli aggregati, evitando una seconda trascrizione manuale.

\section{Run storiche}

\begin{limite}
Le run precedenti all'unificazione dello schema vanno marcate esplicitamente in
questa tabella e nelle sezioni in cui compaiono. In particolare: le run
dell'esperimento~02 anteriori alla campagna definitiva non sono appaiate (dataset
diversi fra la variante classica e quella quantistica) e usano uno schema privo
dei campi \code{test\_accuracy} e \code{dataset}; non sostengono alcuna
conclusione principale.
\end{limite}
